We simulate and analyse the sandpile model for a range of parameter choices and for large measurement counts $n$, as described in the previous chapter. All simulated parameter configurations are summarised in Table~\ref{tab:model_configurations}. For ease of reference, each configuration is assigned a letter label, which is used throughout the remainder of this work.

\begin{table}[h!]
\centering
\begin{tabular}{lccccc}
\hline
config & $d$ & $N$ & boundary condition & perturbation & $n$ \\
\hline
A & 2 & 40 & open   & nonconservative & $10^7$ \\
B & 2 & 40 & open   & conservative    & $10^7$ \\
C & 2 & 40 & closed & nonconservative & $10^7$ \\
D & 3 & 20 & open   & nonconservative & $10^7$ \\
E & 3 & 20 & open   & conservative    & $10^7$ \\
F & 3 & 20 & closed & nonconservative & $10^7$ \\
G & 4 & 20 & open   & nonconservative & $10^7$ \\
H & 4 & 20 & open   & conservative    & $10^7$ \\
I & 4 & 20 & closed & nonconservative & $10^7$ \\
\hline
\end{tabular}
\caption{Simulated model configurations and measurement counts used in this work.}
\label{tab:model_configurations}
\end{table}


Before measurements can be taken, each system must undergo a burn-in phase in order to reach a stationary state. The burn-in procedure is applied as described in the previous section and is monitored using the mean slope $\langle z\rangle(\tau)$ as a function of time $\tau$.
\begin{figure}[h!]
  \centering
  \includegraphics[width=0.8\linewidth]{../../figures/burn_in_plots/N40d2_open_nonconservative_zero_burn_in.png}
  \caption{Burn-in of the mean slope $\langle z\rangle(\tau)$ for system A. The system is initialized with $z=0$ at all sites.}
  \label{fig:N40d2_open_nonconservative_zero_burn_in}
\end{figure}

\begin{figure}[h!]
  \centering
  \includegraphics[width=0.8\linewidth]{../../figures/burn_in_plots/N40d2_open_conservative_zero_burn_in.png}
  \caption{Burn-in of the mean slope $\langle z\rangle(\tau)$ for system B. The system is initialized with $z=0$ at all sites.}
  \label{fig:N40d2_open_conservative_zero_burn_in}
\end{figure}

\begin{figure}[h!]
  \centering
  \includegraphics[width=0.8\linewidth]{../../figures/burn_in_plots/N40d2_open_conservative_random_burn_in.png}
  \caption{Burn-in of the mean slope $\langle z\rangle(\tau)$ for system B. The system is initialized with $z \sim U(0,z_c)$ at all sites.}
  \label{fig:N40d2_open_conservative_random_burn_in}
\end{figure}





Figures~\ref{fig:N40d2_open_nonconservative_zero_burn_in} and~\ref{fig:N40d2_open_conservative_zero_burn_in} show $\langle z\rangle(\tau)$ for configurations A and B when the system is initialized with $z=0$ at all sites. In both cases the approach to a stationary plateau is clearly visible. However, the conservative perturbation scheme (configuration B) requires a substantially longer transient. A plausible explanation is that, under conservative driving, the system effectively gains energy only through additions that occur at the boundary. This slows down the process of reaching the stationary state.

To reduce the burn-in time, an alternative initialization was tested in which each site is assigned an independent random height in the interval $[0,z_c]$, i.e.\ $z \sim U(0,z_c)$. The resulting burn-in curve for configuration B is shown in Figure~\ref{fig:N40d2_open_conservative_random_burn_in}. Compared to the zero initialization, the system reaches the stationary regime significantly faster. For configuration A, the burn-in time does not change, as it is already minimal for the algorithm inroduced in the methods section.

Based on this comparison, random initialization is used for all subsequent simulations and measurements. Burn-in times for all model configurations were then determined using this initialization scheme; the corresponding values are reported in Table~\ref{tab:burn_in_times}.



To extract scaling exponents and the corresponding power-law prefactors, probability distribution functions (PDFs) are computed as described in the previous chapter and evaluated on log-log axes. In such a representation, power-law behaviour appears approximately linear, with the slope corresponding to the scaling exponent. The exponent is therefore estimated by linear regression in the scaling regime. Deviations from linearity are expected both at small values (discreteness and microscopic effects) and at large values (finite-size effects), as discussed in the Theory section, so the regression range must be selected with care.

To determine a suitable fitting window in a reproducible way, a rolling-window procedure is employed. A linear regression is performed on each window, and for the $i$-th window the slope $m_i$, its standard error $\delta m_i$, and the coefficient of determination $r_i^2$ are recorded. Adjacent windows are then compared by forming the slope difference $d_i$ and its propagated uncertainty $\delta d_i$,
\begin{equation}
    d_i = \abs{m_i-m_{i+1}}, \qquad
    \delta d_i = \sqrt{(\delta m_i)^2+(\delta m_{i+1})^2}.
\end{equation}
A window pair is classified as stable if the slope variation is statistically compatible with zero, $d_i < k\,\delta d_i$, and if both fits satisfy $r_i^2 \ge r_{\mathrm{thresh}}$ and $r_{i+1}^2 \ge r_{\mathrm{thresh}}$, where $k$ is a chosen deviation factor and $r_{\mathrm{thresh}}$ is a minimum fit-quality threshold. The final fitting interval is obtained by selecting the longest contiguous stable region and using its corresponding lower and upper boundaries for the exponent regression. At this point it is important to note that the procedure described above is heuristic and does not yield an exact or unique scaling window. The parameters $k$ and $r_{\mathrm{thresh}}$, as well as the window size, are tuned manually and can vary between model configurations.



